% Рисунок с доски: диаграмма деформирования пластического материала
% с площадкой текучести.
\begin{figure}[htbp]
\centering
\begin{tikzpicture}[>={Latex[length=4pt]}]
  \draw[->] (0,0) -- (5.5,0) node[anchor=north east] {$\dev$};
  \draw[->] (0,0) -- (0,3.0) node[anchor=south east] {$\Sdev$};

  \draw[very thick, blue!60!black] (0,0) -- (1.6,2.0) -- (4.8,2.0);
  \draw[dashed, gray] (0,2.0) -- (1.6,2.0);
  \node[anchor=east, font=\footnotesize] at (-0.05,2.0) {$\tau_s$};

  \node[font=\footnotesize, anchor=south] at (3.4,2.1) {пластичность};
  \node[font=\footnotesize, anchor=west, align=left] at (1.0,0.9)
    {упругость};
\end{tikzpicture}
\caption{Пластический материал: пока $N(\Sdev) < \tau_s$, деформация не растёт
($\dt{\dev} = 0$); при $N(\Sdev) = \tau_s$ выполняется $\Sdev = \lambda\dt{\dev}$,
где множитель $\lambda \ge 0$ подстраивается так, чтобы напряжение не росло.}
\label{fig:plastic-diagram}
\end{figure}
